%%%%

%%%   Самарский государственный технический университет

%%%   Кафедра прикладной математики и информатики

%%%

%%%   Преамбула для статьи в журнал "Вестник Самарского государственного

%%%   технического университета. Серия: «Физико-математические науки»"

%%%   Ориентирована под MikTEx 2.4 for Win32

%%%

%%%   Хотя сам журнал собирается под GNU/Linux на дистрибутиве TeXLive



\documentclass[11pt,a4paper,twoside]{article}



\usepackage[cp1251]{inputenc}

\usepackage[T2A]{fontenc}

\usepackage[english,russian]{babel}

\usepackage{samgtubib}

\usepackage[dvips]{graphicx}

\usepackage[multiple]{footmisc}



\usepackage{samgtu}

\renewcommand{\VestnikYear}{2009} % Год издания Вестника

\renewcommand{\VestnikNumber}{2\,(19)} % Номер Вестника

\renewcommand{\VestnikMonth}{Ноябрь} % Месяц выхода Вестника

\renewcommand{\VestnikData}{01.11.09} % Дата подписания Вестника в печать

\renewcommand{\VestnikUPL}{26,64} % Усл. печ. л.

\renewcommand{\VestnikUIL}{25,73} % Уч.-изд. л.

\renewcommand{\VestnikRegNumber}{479/09} % Регистрационный номер

\renewcommand{\VestnikOrder}{994} % Номер заказа






%
%\begin{document}

\renewcommand{\firstpage}{40}
\renewcommand{\lastpage}{45}

\renewcommand{\baselinestretch}{0.95}

\udk{519.633}
\msc{35K10}

\title{Об одной нелокальной краевой задаче для нагруженного дифференциального уравнения со~знакопеременной характеристической формой}
{Об одной нелокальной краевой задаче для нагруженного дифференциального уравнения \dots}

\author{А.\,А.~Токова}{Т\,о\,к\,о\,в\,а~А.\,А.}

\institute{Северо-Кавказская государственная гуманитарно-технологическая академия, \\
369000, Черкесск, ул.~Ставропольская, 36.}

\email{E-mail}{tokova-aa@yandex.ru}

\otherinf{\fullauthor{Алла Аскербиевна Токова} \regalia{к.ф.-м.н.}\position{доцент} \dept{каф. математики}}


\keywords{нагруженное дифференциальное уравнение, нелокальная краевая задача, регулярное решение, существование, единственность}


\workabstract{Решена нелокальная краевая задача для нагруженного уравнения параболического типа со~знакопеременной
характеристической формой. 
Построено общее представление решения. 
Доказаны теорема об общем представлении решения и теорема существования и~единственности решения краевой задачи в~прямоугольной области.}

\engtitle{On a~non-local boundary value problem for~loaded differential equation with characteristic form of~variable sign}

\engauthor{A.\,A.~Tokova}{T\,o\,k\,o\,v\,a~A.\,A.}

\otherenginfm{\fullauthor{Alla A. Tokova}  \regalia{Ph.\,D. (Phys. \& Math.)} \position{Associate Professor} \dept{Dept. of Mathematics}}

\enginstitute{North-Caucasian State Humanitarian-Technological Academy, \\
36,  Stavropol'skaya, Cherkessk,  369000, Russia.  }

\engabstract{Non-local boundary value problem for loaded equation of parabolic type with the sign-variable characteristic form is solved. Common representation of solution is constructed. The theorems of common representation, existence and uniqueness of solution for boundary value problems set in rectangular domain are proved.}

\engkeywords{loaded differential equation, regular solution, non-local boundary value problem, uniqueness, occurence}

\date{08/XI/2010}
\revisiondate{30/V/2011}



\maketitle


В области $\Omega=\{(x,y): -1<x<1,\,0<y<T \}$ рассмотрим нагруженное уравнение
\begin{equation}
\label{tokova:eq1}
L_1 u=L_2 \delta u,
\end{equation}
где
\[
L_1 \equiv x \frac{\partial^2}{\partial x^2} + \lambda^2, \quad L_2 \equiv a \frac{\partial}{\partial y}.
\]
Здесь
$(\delta u )(y)$ "--- среднее значение функции $u(x,\,y)$ по переменной $x$ на сегменте 
$[-1,\,1]$:
$$
(\delta u )(y)=\frac{1}{2}\int _{-1}^{1} u(x,\,y)dx,
$$
$\lambda=\lambda(y)$, $a=a(y)$ "--- заданные непрерывные функции.

Уравнение (\ref{tokova:eq1}) в области $\Omega$ является нагруженным уравнением параболического типа
со знакопеременной характеристической формой  $\Theta(x,y;\,\xi)=x \xi^{2}$.

В работе \cite{nakhushev-du82} А.\,М.~Нахушев предожил метод приближенного решения
краевых задач для дифференциальных
уравнений, основанный на редукции к нагруженным уравнениям
(см.~\cite{nakhushev-du76}).
В связи с этим вызывает интерес постановка и исследование краевых задач
для уравнений вида (\ref{tokova:eq1}).
В даной работе строится общее представление решений уравнения (\ref{tokova:eq1}) и решается
нелокальная краевая задача.
Первая краевая задача для уравнения вида (\ref{tokova:eq1})
рассматривалась в работе \cite{tok5}.

Обозначим через $\Omega^{+}=\Omega \cap \{x>0\}$ и $\Omega^{-}=\Omega \cap\{x<0\}$. Регулярным решением уравнения (\ref{tokova:eq1}) в
области $\Omega$ будем
называть решение $u=u(x,\,y)$ такое, что $u\in C(\bar\Omega)$, $u_{xx} \in C(\Omega^{+})\cap C(\Omega^{-})$, $(\delta u)(y) \in C
[0,\,T] \cap C^1(0,\,T)$ и выполнено соотношение
\[
\lim_{\varepsilon \rightarrow 0} [u_x(\varepsilon,\,y)-u_x(-\varepsilon,\,y)]=0,\quad y\in(0,\,T).
\]

\smallskip

\Section{Представление решения} Будем считать, что $\lambda(y)$, $a(y)>0$, $y\in[0,\,T]$. Пусть $u=u(x,\,y)$ "--- регулярное решение
уравнения (\ref{tokova:eq1}). Обозначим 
\begin{equation}
\label{tokova:eq2}
v(x,\,y)=u(x,\,y)-\frac{a(y)}{\lambda^2(y)}\frac{\partial}{\partial y}(\delta u)(y).
\end{equation}
Тогда функция $v(x,\,y)$ является решением уравнения
\[
x v_{xx}(x,\,y)+\lambda^2(y) v(x,\,y)=0
\]
и в областях $\Omega^{+}$ и $\Omega^{-}$ представляется в виде
\begin{gather}
\label{tokova:eq3}
v(x,\,y) = \sqrt x \left[A(y)I_1 (2\lambda \sqrt{x})+B(y)K_1 (2\lambda \sqrt{x})\right],\quad (x,\,y)\in\Omega^{+},\\
\label{tokova:eq4}
v(x,\,y)=\sqrt{-x}\left[C(y)J_1(2\lambda \sqrt{-x})+D(y)Y_1 (2\lambda \sqrt{-x})\right],\quad (x,\,y)\in\Omega^{-},
\end{gather}
где $I_\nu(z)$, $K_\nu(z)$ "--- модифицированные функции Бесселя 1-го и 2-го рода;
$J_\nu(z)$, $Y_\nu(z)$ "--- функции Бесселя 1"=го и 2"=го рода.

Учитывая формулы дифференцирования функций Бесселя
\begin{gather}
\label{tokova:eq5}
\frac{d}{dz}\left[z^\nu I_\nu(z)\right]=z^\nu I_{\nu-1}(z), \quad
\frac{d}{dz}\left[z^\nu J_\nu(z)\right]=z^\nu J_{\nu-1}(z), \\
\label{tokova:eq6}
\frac{d}{dz}\left[z^\nu K_\nu(z)\right]=-z^\nu
K_{\nu-1}(z), \quad \frac{d}{dz}\left[z^\nu Y_\nu(z)\right]=z^\nu Y_{\nu-1}(z),
\end{gather}
получим
\begin{gather}
\label{tokova:eq7}
v_x (x,\,y)=\lambda \left[A(y)I_0(2 \lambda \sqrt{x})- B(y)K_0(2 \lambda \sqrt{x}) \right],\quad (x,\,y)\in\Omega^{+},\\
\label{tokova:eq8}
v_x(x,\,y)=-\lambda \left[C(y)J_0 (2\lambda \sqrt{-x})  + D(y)Y_0 (2 \lambda \sqrt{-x})\right],\quad (x,\,y)\in\Omega^{-}.
\end{gather}
Из формул \eqref{tokova:eq3}, \eqref{tokova:eq4}, \eqref{tokova:eq7} и \eqref{tokova:eq8}, учитывая асимптотические формулы
\begin{gather*}
I_0(z)=1+{\sf O}(z), \quad K_0(z)=\Psi(1)-\ln\frac{z}{2}+{\sf O}(z\ln z),\\
J_0(z)=1+{\sf O}\left(\frac{1}{z}\right), \quad Y_0(z)=-\frac 2\pi\left[\Psi(1)-\ln\frac{z}{2}\right]+{\sf O}(z\ln z),
\end{gather*}
в силу соотношений
\[
v(x,y)\in C(\overline\Omega),\quad\lim \limits_{\varepsilon \rightarrow 0}
[v_x(\varepsilon,\,y)-v_x(-\varepsilon,\,y)]=0
\]
находим
\begin{equation}
\label{tokova:eq9}
D(y)=-\frac {\pi}{2}B(y),\quad C(y)=-A(y).
\end{equation}

Обозначив 
\begin{equation}
\label{tokova:eq10}
\xi(y)=v(-1,\,y),\quad \eta(y)=v(1,\,y),
\end{equation}
из \eqref{tokova:eq3} и \eqref{tokova:eq4} с учётом \eqref{tokova:eq9} получаем
\begin{gather*}
-A(y)J_1(2\lambda) - \frac{\pi}{2} B(y)Y_1 (2 \lambda)=\xi(y),\quad 
A(y)I_1 (2\lambda)+ B(y)K_1(2 \lambda)=\eta(y).
\end{gather*}
Пусть
\[
\Delta=\frac{\pi}{2}I_1(2\lambda)Y_1 (2\lambda)-J_1(2 \lambda)K_1 (2\lambda),
\]
и $\Delta \neq 0$. Тогда
\begin{equation}
\label{tokova:eq11}
\begin{array}{l}
\displaystyle A(y)=\frac{1}{\Delta} \left[\xi(y)K_1(2 \lambda)+\frac{\pi}{2}\eta(y)Y_1 (2 \lambda)\right],\\ \vspace{2mm}
\displaystyle B(y)=-\frac{1}{\Delta}\left[\xi(y)I_1(2 \lambda)+\eta(y)J_1(2 \lambda)\right].
\end{array}
\end{equation}
Отсюда с учётом \eqref{tokova:eq3}, \eqref{tokova:eq4} и \eqref{tokova:eq9} получаем
\begin{multline}
\label{tokova:eq12}
v(x,\,y)=\frac{\xi(y)\sqrt x}{\Delta} \left[I_1 (2 \lambda \sqrt{x})K_1 (2 \lambda) - K_1 (2 \lambda \sqrt{x}) I_1 (2 \lambda)\right]+ \\
+\frac{\eta(y)\sqrt x}{\Delta}\left[\frac{\pi}{2} I_1 (2\lambda \sqrt{ x}) Y_1 (2 \lambda) - K_1 (2 \lambda \sqrt{x})J_1 (2 \lambda)\right],
\quad (x,\,y) \in\Omega^{+},
\end{multline}
\begin{multline}
\label{tokova:eq13}
v(x,\,y)=\frac{\xi(y)\sqrt {-x}}{\Delta} \left[-J_1(2 \lambda \sqrt{-x})K_1(2 \lambda)+\frac{\pi}{2} Y_1(2 \lambda \sqrt{-x})I_1(2 \lambda)\right]+ \\
+\frac{\eta(y)\sqrt {-x}}{\Delta} \left[-\frac{\pi}{2} J_1 (2 \lambda \sqrt{-x}) Y_1(2 \lambda)+\frac{\pi}{2} Y_1(2 \lambda \sqrt{-x})J_1(2 \lambda)\right],
\, (x,\,y)\in\Omega^{-}.
\end{multline}

Найдём $(\delta v)(y)$. Из \eqref{tokova:eq12}, \eqref{tokova:eq13} получаем
\begin{multline}
\label{tokova:eq14}
2 (\delta v)(y)=\int _{-1}^1 v(x,\,y)dx= \\ = A(y)\left[\int _0^1
\sqrt x I_1(2\lambda \sqrt x) dx -\int _{-1}^0 \sqrt {-x} J_1(2 \lambda \sqrt{-x}) dx \right]+\\
 +B(y)\left[\int _0^1 \sqrt x K_1(2 \lambda \sqrt{x}) dx- \frac{\pi}{ 2}\int _{-1}^0 \sqrt {-x} Y_1 (2 \lambda \sqrt{-x}) dx\right].
\end{multline}
Вычислим интегралы в \eqref{tokova:eq14}. 
Воспользовавшись формулами \eqref{tokova:eq5} и \eqref{tokova:eq6}, сделав замену $x=2\lambda \sqrt t$ в первом и третьем интегралах и замену $x=-2\lambda \sqrt t$ во втором и четвёртом, получаем
\begin{gather*}
\int _0^1 \sqrt x I_1 (2 \lambda \sqrt{x}) dx=\frac{1}{4\lambda^3} \int _0^{2\lambda} t^2 I_1(t)dt= \frac{1}{4\lambda^3}
\int _0^{2\lambda} \left[t^2 I_2(t)\right]'dt=\frac{1}{\lambda} I_2 (2 \lambda),\\
\int _{-1}^0 \sqrt {-x} J_1 (2 \lambda \sqrt{-x}) dx=\frac{1}{4\lambda^3} \int _0^{2\lambda} t^2
J_1(t)dt= \frac{1}{4\lambda^3} \int _0^{2\lambda} \left[t^2 J_2(t)\right]'dt=\frac{1}{\lambda} J_2 (2 \lambda),
\end{gather*}
и, учитывая, что $\lim \limits_{t \rightarrow 0} t^2 K_2(t)=2$
и $\lim \limits_{t \rightarrow 0} t^2 Y_2(t)=-{4}/{\pi}$, имеем
\begin{multline*}
\int _0^1 \sqrt x K_1 (2 \lambda \sqrt x) dx=\frac{1}{4\lambda^3}\int _0^{2\lambda} t^2 K_1(t)dt= \\ =
-\frac{1}{4\lambda^3}\int _0^{2\lambda} \left[t^2 K_2(t)\right]'dt=\frac {1}{2\lambda^3}-\frac{1}{\lambda} K_2(2 \lambda), 
\end{multline*}
\begin{multline*}
\int _{-1}^0 \sqrt{-x} Y_1(2 \lambda \sqrt{-x}) dx=\frac{1}{4\lambda^3}\int _0^{2\lambda} t^2 Y_1(t)dt= \\ =
\frac{1}{4\lambda^3}\int _0^{2\lambda} \left[t^2 Y_2(t)\right]'dt=\frac{1}{\lambda}
Y_2(2 \lambda)+\frac{ 1}{\pi\lambda^3}.
\end{multline*}
Подставляя вычисленные значения в \eqref{tokova:eq14}, получаем
\[
(\delta v)(y)= \frac{A(y)}{2\lambda} \left[I_2(2 \lambda)- J_2(2 \lambda) \right] -  \frac{B(y)}{2\lambda}
\left[K_2(2 \lambda)+ \frac{\pi}{2}Y_2 (2 \lambda)\right],
\]
или с учётом \eqref{tokova:eq11} "---
\begin{equation}
\label{tokova:eq15}
(\delta v)(y)= \xi (y)E(\lambda)+\eta(y)F(y),
\end{equation}
где
\begin{gather*}
E(\lambda)=\frac{1}{2\lambda\Delta} \left(K_1 (2 \lambda)\left[I_2(2 \lambda)- J_2(2 \lambda)\right]+ I_1(2 \lambda)\left[K_2(2 \lambda)+\frac{\pi}{2}
Y_2 (2 \lambda)\right]\right), \\
F(\lambda)= \frac{1}{2\lambda \Delta} \left(\frac{\pi}{2}Y_1(2 \lambda) \left[I_2(2 \lambda)- J_2(2 \lambda)\right] +J_1(2 \lambda)\left[K_2(2 \lambda)+
\frac{\pi}{2} Y_2(2 \lambda)\right]\right).
\end{gather*}
Учитывая, что в силу \eqref{tokova:eq2}
\[
(\delta v)(y)=(\delta u)(y) -\frac{a(y)}{\lambda^2(y)}(\delta u)'(y),
\]
получаем
\begin{equation}
\label{tokova:eq16}
\xi(y)=u(-1,\,y)-\frac{a(y)}{\lambda^2(y)}(\delta
u)'(y), \quad \eta(y)=u(1,\,y)-\frac{a(y)}{\lambda^2(y)}(\delta u)'(y).
\end{equation}
Подставляя \eqref{tokova:eq2} и \eqref{tokova:eq16} в \eqref{tokova:eq15} и считая
$E(\lambda)+F(\lambda)\neq 1$, $y\in[0,\,T]$, приходим к соотношению
\[
(\delta u)'(y) +K(y)(\delta u)(y)= M(y),
\]
где
\[
K(y)=\frac{\lambda^2(y)}{a(y)\left[E(\lambda)+F(\lambda)-1\right]}, \quad M(y)=K(y)\left[E(\lambda)u(-1,\,y)+F(\lambda)u(1,\,y)\right].
\]

Таким образом, получаем
\begin{equation}
\label{tokova:eq17}
(\delta u)(y)=(\delta u)(0)w(y,\,0)+\int _0^y M(t)w(y,\,t)dt,
\end{equation}
где
\[
w(y,\,t)=\exp\left(-\int _t^y K(s)ds\right).
\]

Сформулируем доказанное в виде теоремы.

\smallskip

\hypertarget{tok:th:1}{}

\begin{theorem}[1]Пусть выполняются условия{\/\rm:} 
\begin{itemize} 
\item[\hypertarget{tok:1}{\rm 1)}] $u(x,\,y)$ "--- регулярное решение уравнения {\rm(\ref{tokova:eq1});} 
\item[\hypertarget{tok:2}{\rm 2)}] $a(y),$  $\lambda(y)\in C[0,\,T];$ 
\item[\hypertarget{tok:3}{\rm 3)}] $\Delta \neq 0,$ $E(\lambda)+F(\lambda)\neq 1,$ $a(y)>0,$  $\lambda(y)>0,$ $y \in[0,\,T].$ 
\end{itemize}
Тогда $u(x,\,y)$ представимо в виде
\begin{equation}
\label{tokova:eq18}
u(x,\,y)=v(x,\,y)+\frac{a(y)}{\lambda^2(y)}(\delta u)'(y),
\end{equation}
где $v(x,\,y)$ определяется с помощью \eqref{tokova:eq12}{\rm,} \eqref{tokova:eq13}{\rm ,} а $(\delta u)(y)$ "--- с помощью 
\eqref{tokova:eq17}\rm.
 \end{theorem}

\smallskip

\Section{Краевая задача}
 Сформулируем краевую задачу для уравнения
\eqref{tokova:eq1}: {\sl найти регулярное в области $\Omega $ решение $u(x,\,y)$ уравнения \eqref{tokova:eq1}, удовлетворяющее условиям
\begin{gather}
\label{tokova:eq19}
u(-1,\,y)=\varphi(y), \quad u(1,\,y)=\psi(y), \quad 0<y<T, \\
\label{tokova:eq20}
\alpha(\delta u)(0)+\beta\alpha(\delta u)(T)=\gamma,
\end{gather}
где $\varphi(y)$ и $\psi(y)$ "--- заданные непрерывные функции; $\alpha$, $\beta$, $\gamma$ "--- постоянные, причём $\alpha^2+\beta^2\neq
0$.}

Пусть
\begin{equation}
\label{tokova:eq21}
\alpha+\beta w(T,\,0)\neq 0.
\end{equation}
Подчиняя \eqref{tokova:eq17} условию  \eqref{tokova:eq20}, получим
\begin{gather}
\notag
\gamma=(\delta u)(0)(\alpha+\beta w(T,\,0))+\beta \int _0^T K(t)w(T,\,t)dt,\\
\label{tokova:eq22}
(\delta u)(0)=\frac{\displaystyle\gamma-\beta\int _0^T
K(t)w(T,\,t)dt}{\alpha+\beta w(T,\,0)}.
\end{gather}

С учётом этого из теоремы 1 следует

\smallskip
\begin{theorem}[2]Пусть $\varphi(y),$ $\psi(y)\in C[0,\,T]$ и выполнены условия 
{\rm \hyperlink{tok:2}{2}), \hyperlink{tok:3}{3})} теоремы~{\rm \hyperlink{tok:th:1}{1}} и условие {\rm \eqref{tokova:eq21}.} Тогда существует единственное решение задачи {\rm \eqref{tokova:eq1}, \eqref{tokova:eq19}} и \eqref{tokova:eq20} $u(x,y),$ которое имеет вид {\rm \eqref{tokova:eq18},} 
при этом
\[
\xi(y)=\varphi(y)-\frac{a(y)}{\lambda^2(y)}(\delta u)'(y), \quad \eta(y)=\psi(y)-\frac{a(y)}{\lambda^2(y)}(\delta u)'(y),
\]
а $(\delta u)(y)$ и $(\delta u)(0)$ определены равенствами \eqref{tokova:eq17} и {\rm \eqref{tokova:eq22}} соответственно\rm.
\end{theorem}

\pagebreak

\begin{thebibliography}{9}

\RBibitem{nakhushev-du82}
\by Нахушев~А.\,М.
\paper Об одном приближенном методе решения краевых задач для дифференциальных уравнений и~его приложения к динамике почвенной влаги и грунтовых вод
\jour Дифференц. уравнения
\yr 1982
\vol 18
\issue 1
\pages 72--81
\misctransl
\by Nakhushev~A.\,M. 
\paper An approximate method for solving boundary value problems for differential equations and its application to the dynamics of ground moisture and ground water
\jour Differents. Uravneniya 
\vol 18
\yr 1982
\issue 1
\pages 72–81

\RBibitem{nakhushev-du76}
\by Нахушев~А.\,М.
\paper О задаче Дарбу для одного вырождающегося нагруженного интегро"=дифференциального уравнения второго порядка
\jour Дифференц. уравнения
\yr 1976
\vol 12
\issue 1
\pages 103--108
\misctransl 
\by Nakhushev~A.\,M.
\paper On the Darboux problem for a certain degenerate se\-cond-or\-der loaded integro-differential equation
\jour Differents. Uravneniya 
\vol 12
\yr 1976
\issue 1
\pages 103–108


\RBibitem{tok5}
\by Токова~А.\,А.
\paper О первой краевой задаче для одного нагруженного дифференциального уравнения второго порядка
\jour Докл. Адыгской (Черкесской) Международной академии наук
\yr 2005
\vol 8
\issue 1
\pages 87--91
\misctransl
\paper On the first boundary value problem for a certain second-order loaded  differential equation
\jour Dokl. Adygskoy (Cherkesskoy) Mezhdunarodnoy akademii nauk
\yr 2005
\vol 8
\issue 1
\pages 87--91

\end{thebibliography}




\noindent
\hskip6.5cm \thedate \par\noindent
\hskip6.5cm\therevdate\par


\makeengtitlem

\renewcommand{\baselinestretch}{0.9}


\newpage

%\end{document}
